\documentclass[12pt]{article}
\usepackage[utf8]{inputenc}
\usepackage[T1]{fontenc}
\usepackage[margin=1in]{geometry}
\usepackage{setspace}
\usepackage{xcolor}
\usepackage[colorlinks=true, urlcolor=blue!60!black, linkcolor=blue!60!black]{hyperref}
\setlength{\parindent}{0em}
\setlength{\parskip}{0.45em}
\renewcommand{\baselinestretch}{1.0}
\begin{document}
\begin{center}
{\Large \textit{New York State Rifle \& Pistol Association, Inc. v. Bruen}, 597 U.S. 1 (2022)\par}
\vspace{0.3em}
{\normalsize Case Brief --- One-Page Edition\par}
\end{center}

\textbf{Memory Jogger.} The public-carry case that replaced means-end scrutiny with text and history (\href{https://www.law.cornell.edu/supremecourt/text/20-843}{597 U.S. 1}).

\textbf{Facts.} New York required ``proper cause,'' a special need for self-protection, for an unrestricted concealed-carry license, and banned open carry. Two law-abiding residents were limited to hunting and target-shooting licenses.

\textbf{Procedural History.} The district court dismissed under \textit{Kachalsky}; the Second Circuit affirmed.

\textbf{Votes.} Reversed, 6--3. Thomas wrote for the Court; Alito, Kavanaugh, and Barrett concurred; Breyer dissented, joined by Sotomayor and Kagan.

\textbf{Holding.} The Second Amendment protects public carry for self-defense. If its text covers conduct, the government must show the regulation fits the Nation's historical tradition; the proper-cause requirement does not.

\textbf{Reasoning.} Interest balancing contradicts \textit{Heller}; history shows only bans on carrying to terrify, surety laws presuming a right to carry, and outlier territorial bans. Kavanaugh noted that objective shall-issue licensing survives. Breyer said the test lets judges ``pick their friends out of history's crowd.''

\textbf{Examples.} \emph{Same-Side:} a discretionary character test. \emph{Opposite-Side:} objective shall-issue licensing. \emph{Fence-Sitter:} a default ban on carry on private property open to the public.

\textbf{Critique.} Originalists and progressives alike say the analogue hunt confuses lower courts.

\textbf{Key Quotation.} ``When the Second Amendment's plain text covers an individual's conduct, the Constitution presumptively protects that conduct. The government must then justify its regulation by demonstrating that it is consistent with the Nation's historical tradition of firearm regulation. Only then may a court conclude that the individual's conduct falls outside the Second Amendment's `unqualified command.' '' --- Justice Thomas, Opinion of the Court, 597 U.S. 1.

\end{document}
